(1) 空间均匀、随时间突变的Maxwell方程满足
\(\partial_t\mathbf B=-\nabla\times\mathbf E\)，
\(\partial_t\mathbf D=\nabla\times\mathbf H\)。
跨无自由脉冲的瞬间积分可知
\[
\mathbf B(0^+)=\mathbf B(0^-),\quad
\mathbf D(0^+)=\mathbf D(0^-).
\]
时间界面仍对空间平移不变，故波数 \(k\) 保持不变；
\(\omega_j=ck/n_j\)，因此
\(\omega_2/\omega_1=n_1/n_2\)。
这与空间界面上频率不变、空间波数可变的折射完全不同。

(2) 定义 \(r=n_1/n_2\)。正频波的磁场满足
\(B_y=n_jE_x/c\)，而使用同样 \(+k\) 的负频波有
\(B_y=-n_jE_x/c\)。
时间界面的 \(D,B\) 连续条件为
\[
\varepsilon_2(E_++E_-)=\varepsilon_1E_i,\qquad
n_2(E_+-E_-)=n_1E_i.
\]
于是
\[
\boxed{E_+/E_i=\frac{r^2+r}{2},\qquad
E_-/E_i=\frac{r^2-r}{2}.}
\]
负频分支虽写同一个 \(e^{ikz}\)，其传播能流方向为 \(-z\)。介电常数的阶跃使电场 \(E=D/\varepsilon\) 本身不必连续，错误地令 \(E,H\) 均连续会得到不同甚至违反Maxwell的结果。

(3) 无磁非色散介质单个正或负频平面波的总平均电磁能密度为
\(U_j=\varepsilon_j|E_j|^2/2\)。
所以
\[
\frac{U_++U_-}{U_i}=\frac{1+r^2}{2},\qquad
\frac{U_+-U_-}{U_i}=r.
\]
第二式用到了两支振幅平方之差
\([(r^2+r)^2-(r^2-r)^2]/4=r^3\)。
由于 \(\omega_2/\omega_1=r\)，立刻有
\[
(U_+-U_-)/\omega_2=U_i/\omega_1.
\]
这是一种带符号的作用量不变量，不能将它误称作总能量守恒；
\(r=1\) 时正频幅度为1、负频幅度为0、能量比为1。

(4) 第一界面后两支分别在持续时间 \(T\) 内积累相位
\(e^{\mp i\vartheta}\)，\(\vartheta=\omega_2T\)。
再使用反向时间界面 \(n_2\to n_1\) 的连续 \(D,B\) 方程，去掉不影响物理结果的公共相位，可得
\[
\boxed{A_{\rm out}=\cos\vartheta-\frac{i(r^2+1)}{2r}\sin\vartheta,\quad
B_{\rm out}=\frac{i(r^2-1)}{2r}\sin\vartheta.}
\]
精确恒等式为
\(|A_{\rm out}|^2-|B_{\rm out}|^2=1\)。初末介质相同，故输出负频能量比为
\[
U_{\rm back}/U_i=|B_{\rm out}|^2
=\frac{(r^2-1)^2}{4r^2}\sin^2(\omega_2T).
\]
当 \(r=1\) 或 \(\omega_2T=j\pi\) 时没有净时间反射。总电磁能比
\(U_{\rm out}/U_i=|A|^2+|B|^2=1+2|B|^2\)，新增能量来自外部调节介电常数所做的功，绝非无源恒定介质的放大。
